%! Modeling with Quadratics — Unit 3, Lesson 3.7 — Homework (Answer Key)
\documentclass[10pt]{article}
\usepackage{algebra2-article}
\usepackage{algebra2-key}   % pulls in -boxes; do NOT also load -boxes
\newcommand{\TallMath}[1]{$\displaystyle #1\rule[-1.4em]{0pt}{3.2em}$}

\begin{document}

\pageheader{Unit 3, Lesson 3.7}{Homework --- Answer Key}
\namedateperiod

\vspace{0.10in}

\begin{notesbox}{Practice}
Show your work. \textbf{Build the model, pick the feature the question wants, and answer in context with units.}

\begin{enumerate}[leftmargin=1.9em, itemsep=9pt]
  \item \textbf{Projectile.} A ball is thrown up at $32$ ft/s from a $128$-ft ledge: $h(t)=-16t^2+32t+128$.
    \begin{enumerate}[label=(\alph*), leftmargin=1.9em, itemsep=3pt]
      \item Starting height: \ans{$h(0)=128$ ft}
      \item Maximum height and the time it occurs: \ans{$144$ ft at $t=1$ s}
      \item When it hits the ground: \ans{$t=4$ s (reject $-2$)}
    \end{enumerate}

  \item \textbf{Maximum area (three sides).} A rectangular pen against a barn wall uses $48$ ft of fence on three
        sides ($x=$ width). Write the area model, then find the width, the maximum area, and the dimensions.
    \[ A(x)=\ans{$-2x^2+48x$} \qquad x=\ans{$12$ ft} \qquad \text{max area}=\ans{$288$ ft$^2$ ($12\times24$)} \]

  \item \textbf{Maximum revenue.} A gym has $500$ members paying \$20/month. A study says each \$1 increase loses
        $10$ members ($x=$ number of \$1 increases). Write $R(x)$, then find the price and the maximum revenue.
    \[ R(x)=\ans{$-10x^2+300x+10000$} \qquad \text{price}=\ans{$\$35$} \qquad \text{max }R=\ans{$\$12{,}250$} \]

  \item \textbf{Interpret the features.} A company's monthly revenue is $R(x)=-20x^2+40x+1600$ (dollars, $x=$ \$1
        price increases). State what each feature means \emph{in context}:
    \begin{enumerate}[label=(\alph*), leftmargin=1.9em, itemsep=3pt]
      \item the $y$-intercept $(0,1600)$: \ans{with no price increase, revenue is \$1600}
      \item the vertex $(1,1620)$: \ans{one \$1 increase gives the maximum revenue, \$1620}
    \end{enumerate}

  \item \textbf{Find the error.} Dana models a pen against a wall ($40$ ft of fence, three sides) as $A(x)=x(40-x)$
        and gets a maximum at $x=20$. What did Dana set up wrong, and what is the correct maximum area?
        \ansline{The barn wall replaces the fourth side, so only three sides are fenced: the length is $40-2x$, not $40-x$. The correct model is $A(x)=x(40-2x)=-2x^2+40x$, with maximum at $x=10$: area $200$ ft$^2$ ($10\times20$).}
\end{enumerate}
\end{notesbox}

\begin{extensionbox}
\textbf{1. Same height (systems tie-back).} Two drones: $h_A(t)=-16t^2+64t$ and $h_B(t)=32t$. Other than at launch,
when are they at the same height, and what height is that?
\[ t=\ans{$2$ s} \qquad \text{height}=\ans{$64$ ft} \]

\textbf{2. Reasonable domain.} For the pen in \#2, state the feasible domain of $A(x)$ and explain, in one sentence,
why a width outside it makes no sense.
\ansline{$0<x<24$. A width of $0$ or less is not a real pen, and $x\ge24$ would need $48-2x\le0$ ft for the length --- there's not enough fence.}
\end{extensionbox}

\begin{spiralbox}
\textbf{That's a wrap on Unit 3!} You can now graph a parabola from any of its three forms, factor and solve by every
method (square roots, completing the square, the quadratic formula), work with complex numbers, solve systems, and
--- today --- \emph{build} a quadratic model and read the answer off its vertex and zeros. \textbf{Coming next: the
Unit 3 Review \& Test}, which pulls all of it together. Revisit your warm-ups and exit tickets as you study.
\end{spiralbox}

\begin{teachernote}
Common slips: \#1(b) reporting $t=1$ (the \emph{when}) as the height instead of $144$ ft, or keeping $t=-2$ in (c);
\#2 using $48-x$ instead of $48-2x$ (the headline trap); \#3 setting revenue as price\,$\times$\,\emph{fixed} $500$
members rather than the changing $500-10x$; \#4 is pure interpretation --- credit answers that tie $(0,1600)$ to ``no
increase'' and $(1,1620)$ to ``one \$1 increase, the maximum.'' \#5 is the lesson's headline misconception (three
sides $\Rightarrow$ $40-2x$); the corrected max area is $200$ ft$^2$. Extension \#1 reuses Lesson 3.6 substitution
(the drones meet at $t=2$, both $64$ ft); \#2 is an honest feasible-domain statement.
\end{teachernote}

\end{document}
